% Reviewed translation: PDF pages 236--244 / printed pages 222--230.
% The original scan is authoritative; mathematical tokens were checked against it.
\MathBookSourcePage{236}
\section{关于 "流函数--涡量--压力" 格式的进一步论题}
\label{sec:incompressible-further-sfvp}

本节进一步发展二维 Stokes 问题的 "流函数--涡量--压力" 方法.
特别地, 它在一般情形下为 $\psi_h$, $\omega_h$ 和 $p_h$
导出更好的误差估计. 使用特殊网格还可进一步改进这些估计.

\subsection{误差分析的改进}
\label{subsec:incompressible-further-error-refinement}

我们处于 Subsection~\ref{subsec:incompressible-sfvp-mixed-approximation}
的情形中. 除非另有说明, 使用相同的记号和
\eqref{eq:incompressible-sfvp-standard-fe-spaces}
所定义的具体空间 $\Theta_h=M_h$, $\Phi_h$ 和 $Q_h$.
如该小节末尾所述, Theorem~\ref{thm:incompressible-sfvp-standard-fe-error}
给出的误差界相当差, 因为对 $\widetilde V_h$ 中逼近误差的估计过于粗糙.
现在集中研究这一估计. 下一个 Lemma 在抽象情形下建立
\eqref{eq:primitive-discrete-affine-best-approximation} 的类似结果.

\MathBookSourcePage{237}
\begin{Lemma}\label{lem:incompressible-further-kernel-bound}
若 $M_h\subset\Theta_h$, 则对所有
$\mathbf v=(\operatorname{curl}\phi,\theta)\in\widetilde V$, 有
\begin{equation}\label{eq:incompressible-further-kernel-bound}
\begin{split}
\inf_{\mathbf w_h\in\widetilde V_h}|\mathbf v-\mathbf w_h|
\leq
\inf_{\mathbf v_h=(\operatorname{curl}\phi_h,\theta_h)
\in\widetilde X_h}
\biggl[2|\mathbf v-\mathbf v_h|
+\sup_{\mu_h\in M_h}
\frac{|(\operatorname{curl}(\phi-\phi_h),
\operatorname{curl}\mu_h)|}{\|\mu_h\|_{0,\Omega}}
\biggr].
\end{split}
\end{equation}
将 \eqref{eq:incompressible-further-kernel-bound} 右端的范数
$|\cdot|$ 换成 $\|\cdot\|_{\widetilde X}$ 后,
同一界也适用于
$\inf_{\mathbf w_h\in\widetilde V_h}
\|\mathbf v-\mathbf w_h\|_{\widetilde X}$.
\end{Lemma}

\begin{Proof}
令 $\mathbf v_h=(\operatorname{curl}\phi_h,\theta_h)$
是 $\widetilde X_h$ 的任意元素, 并由
\[
(\tau_h,\mu_h)
=(\theta-\theta_h,\mu_h)
-(\operatorname{curl}(\phi-\phi_h),\operatorname{curl}\mu_h)
\qquad\forall\mu_h\in M_h
\]
定义 $\tau_h\in M_h$. 则
\begin{equation}\label{eq:incompressible-further-tau-bound}
\|\tau_h\|_{0,\Omega}
\leq\|\theta-\theta_h\|_{0,\Omega}
+\sup_{\mu_h\in M_h}
\frac{|(\operatorname{curl}(\phi-\phi_h),
\operatorname{curl}\mu_h)|}{\|\mu_h\|_{0,\Omega}}.
\end{equation}
此外, 由于 $M_h\subset\Theta_h$,
$\mathbf z_h=(0,\tau_h)$ 属于 $\widetilde X_h$.
置 $\mathbf w_h=\mathbf v_h+\mathbf z_h$.
由于 $\mathbf v\in\widetilde V$, 有 $\mathbf w_h\in\widetilde V_h$.
于是
\[
|\mathbf v-\mathbf w_h|
\leq|\mathbf v-\mathbf v_h|+\|\tau_h\|_{0,\Omega},
\]
以及
\[
\|\mathbf v-\mathbf w_h\|_{\widetilde X}
\leq\|\mathbf v-\mathbf v_h\|_{\widetilde X}
+\|\tau_h\|_{0,\Omega}.
\]
\eqref{eq:incompressible-further-kernel-bound}
由 \eqref{eq:incompressible-further-tau-bound}
以及 $\mathbf v_h$ 的任意性得到.
\end{Proof}

考察 \eqref{eq:incompressible-further-kernel-bound} 可见,
困难在于恰当地估计
\[
\sup_{\mu_h\in M_h}
\frac{|(\operatorname{curl}(\phi-\phi_h),
\operatorname{curl}\mu_h)|}{\|\mu_h\|_{0,\Omega}},
\]
因为分子含有 $\mu_h$ 的一阶导数, 而分母没有.
这里不能使用 Section~\ref{sec:incompressible-sfvp} 中
\[
|\mu_h|_{1,\Omega}
\leq\frac Ch\|\mu_h\|_{0,\Omega}
\qquad\forall\mu_h\in M_h
\]
这一简单技巧, 因为它太粗糙.
我们改用 Scholz~\cite{ref:scholz-72} 提出的巧妙论证.
大致说来, 他注意到若取
\[
\phi_h=\mathring P_h\phi,
\]
其中 $\mathring P_h$ 是 $\phi$ 到 $\Phi_h$ 上的投影,
则 $(\operatorname{curl}(\phi-\phi_h),\operatorname{curl}\mu_h)$
化为只在边界单元上取的积分之和.
由于这些单元与 $\Omega$ 中单元总数相比很少,
可为该表达式导出尖锐得多的界. 细节由下一个 Lemma 给出.

\MathBookSourcePage{238}
\begin{Lemma}\label{lem:incompressible-further-boundary-element-bound}
假设 $\Omega$ 是有界凸多边形,
$\mathcal T_h$ 是 $\overline\Omega$ 的一族一致正则剖分.
令整数 $k$ 和实数 $p$ 满足 $1\leq k\leq l$ 和 $2\leq p\leq\infty$.
则存在与 $h$ 无关的常数 $C>0$, 使对所有
$\phi\in W^{k+1,p}(\Omega)\cap H_0^1(\Omega)$, 有
\begin{equation}\label{eq:incompressible-further-boundary-element-bound}
\sup_{\mu_h\in\Theta_h}
\frac{|(\operatorname{curl}(\phi-\mathring P_h\phi),
\operatorname{curl}\mu_h)|}{\|\mu_h\|_{0,\Omega}}
\leq Ch^{k-1/2-1/p}\|\phi\|_{k+1,p,\Omega}.
\end{equation}
\end{Lemma}

\begin{Proof}
首先注意, 因为 $\Omega$ 是凸的, $\Phi_h\subset H_0^1(\Omega)$.
其次, 任取 $\mu_h\in\Theta_h$,
令 $\lambda_h$ 是 $\Phi_h$ 中在所有内部有限元节点上与 $\mu_h$
重合的函数. 特别地, 若 $\Sigma_h$ 表示 $\mathcal T_h$
的边界单元之并, 则
$\operatorname{supp}(\mu_h-\lambda_h)\subset\Sigma_h$.
又因为
\[
(\operatorname{curl}(\phi-\mathring P_h\phi),
\operatorname{curl}\phi_h)=0
\qquad\forall\phi_h\in\Phi_h,
\]
所以
\[
(\operatorname{curl}(\phi-\mathring P_h\phi),
\operatorname{curl}\mu_h)
=\sum_{\kappa\subset\Sigma_h}
\int_\kappa
\operatorname{curl}(\phi-\mathring P_h\phi)
\cdot\operatorname{curl}(\mu_h-\lambda_h)\,dx.
\]
由 Hölder 不等式,
\begin{equation}\label{eq:incompressible-further-holder-boundary}
|(\operatorname{curl}(\phi-\mathring P_h\phi),
\operatorname{curl}\mu_h)|
\leq|\phi-\mathring P_h\phi|_{1,p,\Sigma_h}
|\mu_h-\lambda_h|_{1,q,\Sigma_h},
\end{equation}
其中 $1/p+1/q=1$, $p\geq2$.
附录的逆不等式和
Corollary~\ref{cor:appendix-global-inverse-gradient} 的论证给出
\[
|\mu_h-\lambda_h|_{1,q,\Sigma_h}
\leq C_1h^{-1}(\operatorname{meas}\Sigma_h)^{1/q-1/2}
\|\mu_h-\lambda_h\|_{0,\Sigma_h}.
\]
一方面,
\[
\|\mu_h-\lambda_h\|_{0,\Sigma_h}
\leq C_2\|\mu_h\|_{0,\Sigma_h},
\]
因为 $\mu_h-\lambda_h$ 在所有内部节点上为零,
而在所有边界节点上等于 $\mu_h$.
另一方面,
\[
\operatorname{meas}(\Sigma_h)
\leq C_3h\operatorname{meas}(\Gamma).
\]
因此
\[
|\mu_h-\lambda_h|_{1,q,\Sigma_h}
\leq C_4h^{-1/2-1/p}\|\mu_h\|_{0,\Omega}.
\]
代入 \eqref{eq:incompressible-further-holder-boundary}, 得
\[
|(\operatorname{curl}(\phi-\mathring P_h\phi),
\operatorname{curl}\mu_h)|
\leq C_4h^{-1/2-1/p}
|\phi-\mathring P_h\phi|_{1,p,\Omega}
\|\mu_h\|_{0,\Omega}.
\]
再应用投影估计
\eqref{eq:appendix-ritz-projection-lp-error},
即得到所需结论; 因 $p\geq2$, 它与 $k$ 的值无关.
\end{Proof}

\begin{Remark}\label{rem:incompressible-further-real-order-interpolation}
记
\[
c(\phi,\mu_h)
=(\operatorname{curl}(\phi-\mathring P_h\phi),
\operatorname{curl}\mu_h),
\quad
\forall\phi\in W^{1,p}(\Omega),
\quad\forall\mu_h\in\Theta_h,
\]
并以 $\langle\cdot,\cdot\rangle$ 表示关于 $L^2$ 范数的
$\Theta_h$ 与其对偶 $\Theta_h'$ 之间的对偶配对. 则
\[
c(\phi,\mu_h)=\langle C_h(\phi),\mu_h\rangle,
\]
其中对所有整数 $k\in[1,l]$,
\[
C_h\in\mathcal L(W^{k+1,p}(\Omega)\cap H_0^1(\Omega);\Theta_h'),
\qquad
\|C_h\|\leq Ch^{k-1/2-1/p}.
\]

\MathBookSourcePage{239}
因此在相邻的两个整数 $k$ 值之间插值,
由 Theorem~\ref{thm:lions-peetre-interpolation} 可知,
对所有实数 $k\in[1,l]$, 仍有
$C_h\in\mathcal L(W^{k+1,p}(\Omega)\cap H_0^1(\Omega);\Theta_h')$.
插值公式把
\eqref{eq:incompressible-further-boundary-element-bound}
推广到非整数 $k$:
\begin{equation}\label{eq:incompressible-further-boundary-element-real-order}
\sup_{\mu_h\in\Theta_h}
\frac{|(\operatorname{curl}(\phi-\mathring P_h\phi),
\operatorname{curl}\mu_h)|}{\|\mu_h\|_{0,\Omega}}
\leq Ch^{k-1/2-1/p}\|\phi\|_{k+1,p,\Omega}
\quad\forall\text{ real }k\in[1,l].
\end{equation}
\end{Remark}

利用这两个 Lemma, 可改进 $\widetilde V_h$ 中的逼近结果.

\begin{Lemma}\label{lem:incompressible-further-improved-kernel-approximation}
假设 $\Omega$ 和 $\mathcal T_h$ 如
Lemma~\ref{lem:incompressible-further-boundary-element-bound} 中所述.
对每个实数 $k\in[1,l]$ 和实数 $p\in[2,\infty]$,
存在与 $h$ 无关的常数 $C>0$, 使所有
$\mathbf v=(\operatorname{curl}\phi,\theta=-\Delta\phi)
\in\widetilde V$, 其中
$\phi\in W^{k+1,p}(\Omega)\cap H_0^1(\Omega)$, 满足
\begin{equation}\label{eq:incompressible-further-improved-kernel-approximation}
\left\{
\begin{aligned}
\inf_{\mathbf w_h\in\widetilde V_h}|\mathbf v-\mathbf w_h|
&\leq2\inf_{\theta_h\in\Theta_h}
\|\theta-\theta_h\|_{0,\Omega}
+Ch^{k-1/2-1/p}\|\phi\|_{k+1,p,\Omega},\\
\inf_{\mathbf w_h\in\widetilde V_h}
\|\mathbf v-\mathbf w_h\|_{\widetilde X}
&\leq2\left\{
|\phi-\mathring P_h\phi|_{1,s,\Omega}
+\inf_{\theta_h\in\Theta_h}\|\theta-\theta_h\|_{0,\Omega}
\right\}\\
&\quad+Ch^{k-1/2-1/p}\|\phi\|_{k+1,p,\Omega}.
\end{aligned}
\right.
\end{equation}
当
$\phi\in H^{l+3/2}(\Omega)\cap W^{l+1,\infty}(\Omega)$
(当然 $\phi|_\Gamma=0$) 时, 最佳估计为
\begin{equation}\label{eq:incompressible-further-best-kernel-approximation}
\left\{
\begin{aligned}
\inf_{\mathbf w_h\in\widetilde V_h}|\mathbf v-\mathbf w_h|,\\
\inf_{\mathbf w_h\in\widetilde V_h}
\|\mathbf v-\mathbf w_h\|_{\widetilde X}
\end{aligned}
\right\}
\leq Ch^{l-1/2}
\left(\|\phi\|_{l+3/2,\Omega}
+\|\phi\|_{l+1,\infty,\Omega}\right).
\end{equation}
\end{Lemma}

\begin{Proof}
回顾
\[
\|(\operatorname{curl}\phi,\theta)\|_{\widetilde X}
=|\phi|_{1,s,\Omega}+\|\theta\|_{0,\Omega}.
\]
不等式
\eqref{eq:incompressible-further-improved-kernel-approximation}
由 Lemma~\ref{lem:incompressible-further-kernel-bound},
Lemma~\ref{lem:incompressible-further-boundary-element-bound}
和 Remark~\ref{rem:incompressible-further-real-order-interpolation}
直接得到.

其次, 令
$\phi\in H^{l+3/2}(\Omega)\cap W^{l+1,\infty}(\Omega)$.
在 Lemma~\ref{lem:incompressible-further-boundary-element-bound}
中取 $k=l$, $p=\infty$, 得
\[
\sup_{\mu_h\in\Theta_h}
\frac{|(\operatorname{curl}(\phi-\mathring P_h\phi),
\operatorname{curl}\mu_h)|}{\|\mu_h\|_{0,\Omega}}
\leq C_1h^{l-1/2}\|\phi\|_{l+1,\infty,\Omega}.
\]
此外, 因
$H^{l+3/2}(\Omega)\subset W^{l+1/2,s}(\Omega)$
对所有 $s<\infty$ 成立, 应用
\eqref{eq:appendix-ritz-projection-lp-error}, 得
\[
|\phi-\mathring P_h\phi|_{1,s,\Omega}
\leq C_2h^{l-1/2}\|\phi\|_{l+1/2,s,\Omega}
\leq C_3h^{l-1/2}\|\phi\|_{l+3/2,\Omega}.
\]

\MathBookSourcePage{240}
最后, $\theta=-\Delta\phi\in H^{l-1/2}(\Omega)$,
由附录投影估计 (当 $l\geq3$ 时也可用插值估计),
\[
\inf_{\theta_h\in\Theta_h}\|\theta-\theta_h\|_{0,\Omega}
\leq C_4h^{l-1/2}\|\theta\|_{l-1/2,\Omega}
\leq C_4h^{l-1/2}\|\phi\|_{l+3/2,\Omega}.
\]
因此 \eqref{eq:incompressible-further-best-kernel-approximation}
由 \eqref{eq:incompressible-further-improved-kernel-approximation}
和这三个不等式得到.
\end{Proof}

利用简单而经典的稠密性论证, 立即得到下列推论.

\begin{Corollary}\label{cor:incompressible-further-kernel-density}
在 Lemma~\ref{lem:incompressible-further-boundary-element-bound}
的假设下, 对所有 $\mathbf v\in\widetilde V$, 有
\[
\lim_{h\to0}\inf_{\mathbf w_h\in\widetilde V_h}
\|\mathbf v-\mathbf w_h\|_{\widetilde X}=0.
\]
\end{Corollary}

此外, Lemma~\ref{lem:incompressible-further-improved-kernel-approximation}
还蕴含 "流函数--涡量" 方法的若干误差估计.
先考虑 Stokes 问题的解具有足够正则性的最简单情形.

\begin{Theorem}\label{thm:incompressible-further-smooth-error}
令 $\Omega$ 是有界凸多边形,
$\mathcal T_h$ 是 $\overline\Omega$ 的一族一致正则剖分.
假设 Stokes 问题 \eqref{eq:incompressible-sfvp-stokes}
的解 $(\mathbf u=\operatorname{curl}\psi,p)$ 满足
\[
\psi\in H^{k+3/2}(\Omega)\cap W^{k+1,\infty}(\Omega),
\qquad
p\in H^{k-1/2}(\Omega)\cap L_0^2(\Omega),
\]
其中某个实数 $k\in[3/2,l]$.
则问题 \eqref{eq:incompressible-sfvp-practical-problem}
和 \eqref{eq:incompressible-sfvp-discrete-pressure}
的解
$(\mathbf u_h=(\operatorname{curl}\psi_h,\omega_h),p_h)$ 满足
\begin{equation}\label{eq:incompressible-further-smooth-velocity-error}
\|\mathbf u-\mathbf u_h\|_{\widetilde X}
\leq C_1h^{k-1/2}
\left(\|\psi\|_{k+3/2,\Omega}
+\|\psi\|_{k+1,\infty,\Omega}\right),
\end{equation}
以及
\begin{equation}\label{eq:incompressible-further-smooth-pressure-error}
\|p-p_h\|_{0,\Omega}
\leq C_2h^{k-1/2}
\left(|p|_{k-1/2,\Omega}
+\|\psi\|_{k+3/2,\Omega}
+\|\psi\|_{k+1,\infty,\Omega}\right).
\end{equation}
当 $1\leq k\leq\min(3/2,l)$ 且
$\psi\in W^{k+1,\infty}(\Omega)$,
$\Delta\psi,p\in H^1(\Omega)$ 时, 有
\begin{equation}\label{eq:incompressible-further-low-regularity-velocity-error}
\|\mathbf u-\mathbf u_h\|_{\widetilde X}
\leq C_3\left(
h|\Delta\psi|_{1,\Omega}
+h^{k-1/2}\|\psi\|_{k+1,\infty,\Omega}\right),
\end{equation}
以及
\begin{equation}\label{eq:incompressible-further-low-regularity-pressure-error}
\|p-p_h\|_{0,\Omega}
\leq C_4\left(
h|p|_{1,\Omega}+h|\Delta\psi|_{1,\Omega}
+h^{k-1/2}\|\psi\|_{k+1,\infty,\Omega}\right).
\end{equation}
\end{Theorem}

\begin{Proof}
由 \eqref{eq:incompressible-further-improved-kernel-approximation},
Remark~\ref{rem:incompressible-sfvp-sharper-seminorm-error},
Theorem~\ref{thm:incompressible-sfvp-velocity-vorticity-error}
和 Theorem~\ref{thm:incompressible-sfvp-pressure-error}, 得
\[
|\mathbf u-\mathbf u_h|
\leq C_1\left\{
\|\omega-P_h\omega\|_{0,\Omega}
+h^{k-1/2}\|\psi\|_{k+1,\infty,\Omega}\right\},
\]
\[
\|\mathbf u-\mathbf u_h\|_{\widetilde X}
\leq C_2\left\{
\|\omega-P_h\omega\|_{0,\Omega}
+|\psi-\mathring P_h\psi|_{1,s,\Omega}
+h^{k-1/2}\|\psi\|_{k+1,\infty,\Omega}\right\},
\]
以及
\[
\begin{split}
\|p-p_h\|_{0,\Omega}
\leq C_3\biggl\{&h\inf_{q_h\in Q_h}|p-q_h|_{1,\Omega}
+|\mathbf u-\mathbf u_h|\\
&+\inf_{\theta_h\in\Theta_h}
\left(\|\omega-\theta_h\|_{0,\Omega}
+h|\omega-\theta_h|_{1,\Omega}\right)\biggr\}.
\end{split}
\]
应用附录投影估计, 按 $k$ 的取值分别得到
\eqref{eq:incompressible-further-smooth-velocity-error},
\eqref{eq:incompressible-further-smooth-pressure-error}
或 \eqref{eq:incompressible-further-low-regularity-velocity-error},
\eqref{eq:incompressible-further-low-regularity-pressure-error}.
\end{Proof}

\MathBookSourcePage{241}
当 $l\geq2$ 且
$\psi\in H^{l+3/2}(\Omega)\cap W^{l+1,\infty}(\Omega)$,
$p\in H^{l-1/2}(\Omega)$ 时, 最佳估计为
\[
\|\mathbf u-\mathbf u_h\|_{\widetilde X}
\leq Ch^{l-1/2}
\left(\|\psi\|_{l+3/2,\Omega}
+\|\psi\|_{l+1,\infty,\Omega}\right),
\]
\[
\|p-p_h\|_{0,\Omega}
\leq Ch^{l-1/2}
\left(|p|_{l-1/2,\Omega}
+\|\psi\|_{l+3/2,\Omega}
+\|\psi\|_{l+1,\infty,\Omega}\right).
\]
当 $l=1$, $\psi\in W^{2,\infty}(\Omega)$,
$\Delta\psi,p\in H^1(\Omega)$ 时, 最佳估计为
\[
\|\mathbf u-\mathbf u_h\|_{\widetilde X}
\leq Ch\left(|\Delta\psi|_{1,\Omega}
+h^{-1/2}\|\psi\|_{2,\infty,\Omega}\right),
\]
\[
\|p-p_h\|_{0,\Omega}
\leq Ch\left(|p|_{1,\Omega}+|\Delta\psi|_{1,\Omega}
+h^{-1/2}\|\psi\|_{2,\infty,\Omega}\right).
\]

\begin{Remark}\label{rem:incompressible-further-near-smooth-error}
当 $\psi\in H^{k+2}(\Omega)$ 而不属于
$W^{k+1,\infty}(\Omega)$ 时,
Theorem~\ref{thm:incompressible-further-smooth-error} 的估计几乎仍成立.
更精确地, 对这个 $k$ 和任意 $p\geq2$ 应用
Lemma~\ref{lem:incompressible-further-boundary-element-bound},
并使用 Sobolev Theorem~\ref{thm:sobolev-embedding},
可把 \eqref{eq:incompressible-further-improved-kernel-approximation} 换成
\[
\inf_{\mathbf w_h\in\widetilde V_h}
\|\mathbf v-\mathbf w_h\|_{\widetilde X}
\leq C_ph^{k-1/2-1/p}\|\phi\|_{k+2,\Omega}
\qquad\forall p\geq2.
\]
令 $p$ 趨于无穷, 可见对每个 $\varepsilon>0$,
存在 $C(\varepsilon)>0$, 使
\[
\inf_{\mathbf w_h\in\widetilde V_h}
\|\mathbf v-\mathbf w_h\|_{\widetilde X}
\leq C(\varepsilon)h^{k-1/2-\varepsilon}
\|\phi\|_{k+2,\Omega}.
\]
进而
\[
\|\mathbf u-\mathbf u_h\|_{\widetilde X}
\leq C(\varepsilon)h^{k-1/2-\varepsilon}
\|\psi\|_{k+2,\Omega},
\]
这里对 $k$ 或 $\Delta\psi$ 没有附加限制,
因为此时 $\Delta\psi\in H^k(\Omega)$ 且 $k\geq1$.
\end{Remark}

同样, 标准稠密性论证证明 "流函数--涡量" 方法收敛.

\begin{Corollary}\label{cor:incompressible-further-general-convergence}
令 $\Omega$ 和 $\mathcal T_h$ 如
Theorem~\ref{thm:incompressible-further-smooth-error} 中所述,
且 Stokes 问题 \eqref{eq:incompressible-sfvp-stokes}
的解 $(\mathbf u=\operatorname{curl}\psi,p)$
属于 $H_0^1(\Omega)^2\times L_0^2(\Omega)$.
则问题 \eqref{eq:incompressible-sfvp-practical-problem}
和 \eqref{eq:incompressible-sfvp-discrete-pressure} 的解
$(\mathbf u_h,p_h)$ 满足
\[
\lim_{h\to0}
\left(\|\mathbf u-\mathbf u_h\|_{\widetilde X}
+\|p-p_h\|_{0,\Omega}\right)=0.
\]
\end{Corollary}

在 Chapter~\ref{chap:navier-stokes-theory-approx} 研究 Navier--Stokes 方程时,
会遇到正则性不超过 $L^r$ 的右端项 $\mathbf f$, $1<r\leq2$.
由于假设 $\Omega$ 为凸域, 当右端项属于 $L^r(\Omega)^2$ 时,
Stokes 问题的解 $(\psi,\omega,p)$ 属于
$W^{3,r}(\Omega)\times W^{1,r}(\Omega)\times W^{1,r}(\Omega)$.
此时不能直接用附录投影估计来计算
$\|\omega-P_h\omega\|_{0,\Omega}$.
下面证明相应的逼近结果.

\begin{Lemma}\label{lem:incompressible-further-low-r-projection}
保留 Theorem~\ref{thm:incompressible-further-smooth-error}
关于 $\Omega$ 和 $\mathcal T_h$ 的假设.
令 $1<q\leq2$ 且 $1/p+1/q=1$.
则存在与 $h$ 无关的常数 $C>0$, 使
\begin{equation}\label{eq:incompressible-further-low-r-projection}
\|v-P_hv\|_{0,\Omega}
\leq Ch^{2/p}|\ln h|^\beta|v|_{1,q,\Omega}
\qquad\forall v\in W^{1,q}(\Omega),
\end{equation}
其中当 $l\geq2$ 时 $\beta=0$,
当 $l=1$ 时 $\beta=1-2/p$.
\end{Lemma}

\MathBookSourcePage{242}
\begin{Proof}
令 $R_h\in\mathcal L(H^1(\Omega);\Theta_h)$
为附录定义的局部正则化算子, 并写成
\[
\|v-P_hv\|_{0,\Omega}
\leq\|v-R_hv\|_{0,\Omega}
+\|R_hv-P_hv\|_{0,\Omega}.
\]
由于 $R_hv-P_hv\in\Theta_h$,
Lemma~\ref{lem:appendix-global-inverse-same-order} 蕴含
\[
\|R_hv-P_hv\|_{0,\Omega}
\leq Ch^{1-2/q}\|R_hv-P_hv\|_{0,q,\Omega}.
\]
因此
\[
\begin{split}
\|v-P_hv\|_{0,\Omega}
\leq{}&\|v-R_hv\|_{0,\Omega}\\
&+Ch^{1-2/q}
\left(\|v-R_hv\|_{0,q,\Omega}
+\|v-P_hv\|_{0,q,\Omega}\right).
\end{split}
\]
故 \eqref{eq:incompressible-further-low-r-projection}
是附录局部正则化误差、Jensen 不等式
\eqref{eq:appendix-discrete-jensen}
和投影估计的直接推论; 当 $q<2$ 时带有一个对数因子.
\end{Proof}

该 Lemma 使我们能够推广
Theorem~\ref{thm:incompressible-further-smooth-error}.

\begin{Theorem}\label{thm:incompressible-further-w3alpha-error}
令 $\Omega$ 和 $\mathcal T_h$ 如
Theorem~\ref{thm:incompressible-further-smooth-error} 中所述,
并假设流函数 $\psi\in W^{3,\alpha}(\Omega)$,
压力 $p\in W^{1,\alpha}(\Omega)$, 其中 $\alpha\in[r,2)$.
令 $\beta$ 满足 $1/\alpha+1/\beta=1$.
则
\begin{equation}\label{eq:incompressible-further-w3alpha-velocity-error}
\|\mathbf u-\mathbf u_h\|_{\widetilde X}
\leq C_1
\begin{cases}
h^{2/\beta}\|\psi\|_{3,\alpha,\Omega},&l\geq2,\\
h^{1/\beta}\|\psi\|_{3,\alpha,\Omega},&l=1,
\end{cases}
\end{equation}
并且
\begin{equation}\label{eq:incompressible-further-w3alpha-pressure-error}
\|p-p_h\|_{0,\Omega}
\leq C_2(\|\psi\|_{3,\alpha,\Omega}+|p|_{1,\alpha,\Omega})
\begin{cases}
h^{2/\beta},&l\geq2,\\
h^{1/\beta},&l=1,
\end{cases}
\end{equation}
其中 $C_1,C_2$ 与 $h$, $\psi$, $p$ 无关.
\end{Theorem}

\begin{Proof}
当 $l\geq2$ 时, 可应用 Subsection~\ref{subsec:incompressible-sfvp-mesh-dependent}
的结果. Lemma~\ref{lem:incompressible-sfvp-mesh-zero-approximation}
和 Lemma~\ref{lem:incompressible-sfvp-mesh-two-approximation}
的论证可直接推广为
\[
\inf_{\theta_h\in\Theta_h}\|\omega-\theta_h\|_{0,h}
\leq C_1h^{2/\beta}|\omega|_{1,\alpha,\Omega},
\qquad l\geq1,
\]
以及
\[
\inf_{\phi_h\in\Phi_h}\|\psi-\phi_h\|_{2,h}
\leq C_2h^{2/\beta}|\psi|_{3,\alpha,\Omega},
\qquad l\geq2.
\]
注意 $W^{3,\alpha}(\Omega)\subset H^2(\Omega)$.
代入 \eqref{eq:incompressible-sfvp-mesh-abstract-errors},
得到比 \eqref{eq:incompressible-further-w3alpha-velocity-error}
略尖锐的估计
\begin{equation}\label{eq:incompressible-further-mesh-w3alpha-error}
\|\omega-\omega_h\|_{0,h}
+\|\psi-\psi_h\|_{2,h}
\leq C_3h^{2/\beta}|\psi|_{3,\alpha,\Omega},
\qquad l\geq2.
\end{equation}
而 \eqref{eq:incompressible-sfvp-pressure-error}
和 \eqref{eq:incompressible-further-w3alpha-velocity-error}
立即给出 \eqref{eq:incompressible-further-w3alpha-pressure-error}.

当 $l=1$ 时, 必须使用
\eqref{eq:incompressible-sfvp-quasioptimal-error}
和 \eqref{eq:incompressible-further-improved-kernel-approximation}:
\begin{equation}\label{eq:incompressible-further-linear-element-bound}
\begin{split}
\|\mathbf u-\mathbf u_h\|_{\widetilde X}
\leq C_4\bigl\{&\|\omega-P_h\omega\|_{0,\Omega}
+|\psi-\mathring P_h\psi|_{1,s,\Omega}\\
&+h^{k-1/2-1/p}\|\psi\|_{k+1,p,\Omega}\bigr\}.
\end{split}
\end{equation}

\MathBookSourcePage{243}
一方面, Sobolev Theorem~\ref{thm:sobolev-embedding}
蕴含 $W^{3,\alpha}(\Omega)\subset W^{2,p}(\Omega)$,
其中 $1/p=1/\alpha-1/2>0$.
因此在 \eqref{eq:incompressible-further-linear-element-bound}
最后一项中取此 $p$ 和 $k=1$, 有
\[
h^{1/2-1/p}\|\psi\|_{2,p,\Omega}
\leq C_5h^{1/\beta}\|\psi\|_{3,\alpha,\Omega}.
\]
另一方面, 再次应用该 Theorem 可得
$W^{3,\alpha}(\Omega)\subset W^{1+2/\beta,p}(\Omega)$
对所有 $p>0$ 成立. 因此投影估计特别给出
\[
|\psi-\mathring P_h\psi|_{1,s,\Omega}
\leq C_6h^{2/\beta}\|\psi\|_{3,\alpha,\Omega}.
\]
最后, Lemma~\ref{lem:incompressible-further-low-r-projection} 给出
\[
\|\omega-P_h\omega\|_{0,\Omega}
\leq C_7h^{2/\beta}|\ln h|^{1-2/\beta}
|\omega|_{1,\alpha,\Omega}.
\]
所以 \eqref{eq:incompressible-further-linear-element-bound}
中的主导幂为 $h^{1/\beta}$,
这证明 \eqref{eq:incompressible-further-w3alpha-velocity-error}.
再由 \eqref{eq:incompressible-sfvp-pressure-error}
得到 \eqref{eq:incompressible-further-w3alpha-pressure-error}.
\end{Proof}

\begin{Remark}\label{rem:incompressible-further-alpha-two}
当 $\alpha=2$ 时,
Theorem~\ref{thm:incompressible-further-w3alpha-error}
的结论对 $l=2$ 仍成立, 但因
\eqref{eq:incompressible-further-linear-element-bound}
的最后一项, 对 $l=1$ 不成立.
与 Remark~\ref{rem:incompressible-further-near-smooth-error} 类似,
此时得到
\[
\left.
\begin{aligned}
\|\mathbf u-\mathbf u_h\|_{\widetilde X},\\
\|p-p_h\|_{0,\Omega}
\end{aligned}
\right\}
\leq C(\varepsilon)h^{1/2-\varepsilon}
\begin{cases}
\|\psi\|_{3,\Omega},\\
|p|_{1,\Omega}+\|\psi\|_{3,\Omega},
\end{cases}
\qquad\forall\varepsilon>0.
\]
\end{Remark}

最后, 利用熟悉的对偶论证, 下一个 Theorem
补充 Theorem~\ref{thm:incompressible-sfvp-mesh-h1-error},
并在 $l=1$ 时为 $|\psi-\psi_h|_{1,\Omega}$
建立几乎最优的上界.

\begin{Theorem}\label{thm:incompressible-further-h1-error}
令 $\Omega$ 和 $\mathcal T_h$ 如
Theorem~\ref{thm:incompressible-further-smooth-error} 中所述.
假设 Stokes 问题 \eqref{eq:incompressible-sfvp-stokes}
的解 $\mathbf u=\operatorname{curl}\psi$ 满足
\[
\psi\in H^{k+1}(\Omega),\quad 2\leq k\leq l,
\qquad\text{或}\qquad
\psi\in H^3(\Omega),\quad l=1.
\]
则问题 \eqref{eq:incompressible-sfvp-practical-problem}
的解 $\mathbf u_h=(\operatorname{curl}\psi_h,\omega_h)$ 满足
\begin{equation}\label{eq:incompressible-further-h1-error}
|\psi-\psi_h|_{1,\Omega}
\leq
\begin{cases}
Ch^k|\psi|_{k+1,\Omega},&2\leq k\leq l,\\
C(\varepsilon)h^{1-\varepsilon}\|\psi\|_{3,\Omega},&l=1,
\end{cases}
\end{equation}
其中 $\varepsilon>0$ 任意.
\end{Theorem}

\begin{Proof}
当 $l\geq2$ 时,
\eqref{eq:incompressible-further-h1-error}
已由 Theorem~\ref{thm:incompressible-sfvp-mesh-h1-error} 证明.
当 $l=1$ 时, 使用非常相似的对偶论证.
对每个 $\mathbf g\in L^2(\Omega)^2$, 辅助 Stokes 问题
\[
(\operatorname{curl}\lambda_g,\operatorname{curl}\chi)
=(\mathbf g,\operatorname{curl}\chi)
\qquad\forall\chi\in H_0^1(\Omega),
\]
\[
(\operatorname{curl}\phi_g,\operatorname{curl}\mu)
=(\lambda_g,\mu)
\qquad\forall\mu\in H^1(\Omega)
\]
有唯一解 $\lambda_g\in H^1(\Omega)$,
$\phi_g\in H^3(\Omega)\cap H_0^1(\Omega)$, 且
\begin{equation}\label{eq:incompressible-further-dual-regularity}
\|\phi_g\|_{3,\Omega}+\|\lambda_g\|_{1,\Omega}
\leq C_1\|\mathbf g\|_{0,\Omega}.
\end{equation}

\MathBookSourcePage{244}
此外, 有恒等式
\[
\begin{split}
(\mathbf g,\operatorname{curl}(\psi-\psi_h))
={}&(\operatorname{curl}(\lambda_g-\lambda_h),
\operatorname{curl}(\psi-\psi_h))\\
&+(\omega-\omega_h,\lambda_h-\lambda_g)\\
&+(\operatorname{curl}(\phi_g-\phi_h),
\operatorname{curl}(\omega-\omega_h))
\end{split}
\]
对所有 $\lambda_h\in\Theta_h$, $\phi_h\in\Phi_h$ 成立.
在前两项中可取 $\lambda_h=P_h\lambda_g$, 得
\[
(\operatorname{curl}(\lambda_g-P_h\lambda_g),
\operatorname{curl}(\psi-\chi_h))
+(\omega-\omega_h,P_h\lambda_g-\lambda_g)
\qquad\forall\chi_h\in\Phi_h.
\]
另一方面, 第三项可分成
\[
\begin{split}
(\operatorname{curl}(\phi_g-\phi_h),
\operatorname{curl}(\omega-\omega_h))
={}&(\operatorname{curl}(\phi_g-\phi_h),
\operatorname{curl}(\omega-P_h\omega))\\
&+(\operatorname{curl}(\phi_g-\phi_h),
\operatorname{curl}(P_h\omega-\omega_h)).
\end{split}
\]
但
\[
(\operatorname{curl}\phi_g,
\operatorname{curl}(P_h\omega-\omega_h))
=(\lambda_g,P_h\omega-\omega_h),
\]
且由 \eqref{eq:incompressible-sfvp-vorticity-orthogonality}
和 \eqref{eq:incompressible-sfvp-projected-orthogonality}, 有
\[
(P_h\omega-\omega_h,\theta_h)=0
\quad\forall(\operatorname{curl}\sigma_h,\theta_h)
\in\widetilde V_h,
\]
\[
(\operatorname{curl}(P_h\omega-\omega_h),
\operatorname{curl}\phi_h)=0
\qquad\forall\phi_h\in\Phi_h.
\]
汇总这些等式, 得
\begin{equation}\label{eq:incompressible-further-duality-identity}
\begin{split}
(\mathbf g,\operatorname{curl}(\psi-\psi_h))
={}&(\operatorname{curl}(\lambda_g-P_h\lambda_g),
\operatorname{curl}(\psi-\chi_h))\\
&+(\omega-\omega_h,P_h\lambda_g-\lambda_g)\\
&+(\operatorname{curl}(\phi_g-\phi_h),
\operatorname{curl}(\omega-P_h\omega))\\
&+(\lambda_g-\theta_h,P_h\omega-\omega_h)
\end{split}
\end{equation}
对所有 $\chi_h,\phi_h\in\Phi_h$
和所有 $(\operatorname{curl}\sigma_h,\theta_h)\in\widetilde V_h$ 成立.
由于 $\phi_g$ 和 $\psi$ 至少属于 $H^2(\Omega)$, 取
$\chi_h=I_h\psi$, $\phi_h=I_h\phi_g$.
则
\[
|(\operatorname{curl}(\lambda_g-P_h\lambda_g),
\operatorname{curl}(\psi-I_h\psi))|
\leq C_2h|\lambda_g|_{1,\Omega}|\psi|_{2,\Omega},
\]
\[
|(\omega-\omega_h,P_h\lambda_g-\lambda_g)|
\leq C_3h|\omega|_{1,\Omega}|\lambda_g|_{1,\Omega},
\]
\[
|(\operatorname{curl}(\phi_g-I_h\phi_g),
\operatorname{curl}(\omega-P_h\omega))|
\leq C_4h|\phi_g|_{2,\Omega}|\omega|_{1,\Omega}.
\]
又由 Remark~\ref{rem:incompressible-further-near-smooth-error},
\[
\inf_{(\operatorname{curl}\sigma_h,\theta_h)\in\widetilde V_h}
\|\lambda_g-\theta_h\|_{0,\Omega}
\leq C_5(\varepsilon)h^{1/2-\varepsilon}|\phi_g|_{3,\Omega},
\]
而类似地由 Remark~\ref{rem:incompressible-further-alpha-two},
\[
\|P_h\omega-\omega_h\|_{0,\Omega}
\leq C_6(\varepsilon)h^{1/2-\varepsilon}|\omega|_{1,\Omega}
\qquad\forall\varepsilon>0.
\]
把这些估计代入
\eqref{eq:incompressible-further-duality-identity},
再应用 \eqref{eq:incompressible-further-dual-regularity},
即得到 \eqref{eq:incompressible-further-h1-error}.
\end{Proof}
